% comment.tex -- public comment, FDA docket FDA-2025-N-6075.
% Build:  tectonic comment.tex   (run from reports/docket_comment/)
% Nothing in this file is computed here; every model statement is quoted from the memo.

\documentclass[11pt,letterpaper]{article}
\usepackage[T1]{fontenc}
\usepackage[utf8]{inputenc}
\usepackage[letterpaper,margin=1in]{geometry}
\usepackage{newtxtext}
\usepackage{newtxmath}
\usepackage{microtype}
\usepackage[hidelinks]{hyperref}
\usepackage{titlesec}
\usepackage{enumitem}
\titleformat{\section}{\normalfont\normalsize\bfseries}{\thesection.}{0.6em}{}
\titlespacing*{\section}{0pt}{2.0ex plus 0.4ex}{0.8ex}
\setlength{\parskip}{0.7ex plus 0.2ex}
\setlist[enumerate]{leftmargin=1.4em,itemsep=0.4ex,topsep=0.5ex}
\hypersetup{pdftitle={Comment on FDA-2025-N-6075},pdfauthor={Rikhin Kavuru}}

\begin{document}
\pagestyle{plain}

{\bfseries Comment on Docket No.\ FDA-2025-N-6075, RIN 0910-AI94}\\
Drug Establishment Registration and Drug Listing Requirements for Establishments Engaged in
Distributed Manufacturing and Certain Foreign Establishments\\[0.7ex]
Submitted by Rikhin Kavuru, individual.

\vspace{0.4ex}\hrule\vspace{1.2ex}

\section{What this comment is, and what it is not}

\textbf{Disclosure of interest.} I work on early-stage software for sterile drug manufacturing, a
release-decision support layer, and I have sought pre-seed funding and grant support for that work.
It has no product on the market, no customer, no revenue, and no application pending before the
Agency. A rule that makes distributed manufacturing registration more workable could in principle
enlarge the market that work is directed at, so I have an interest in this rulemaking that is
indirect but real. I am submitting as an individual and not on behalf of that work, and I am stating
the interest here rather than leaving the Agency to discover it, so this comment can be weighed
accordingly.

I am a high school student. Over the past year I built an open, pre-registered engineering and
operations feasibility model of sterile-injectable supply, in order to test whether regionally
distributed manufacturing reduces shortage risk more effectively than safety stock, dual sourcing,
reserved contract capacity, or additional centralized capacity. The protocol, the code, the frozen
source snapshots, and the study's own retractions are public.

I want to be direct about the model's standing before making any recommendation. Its numeric inputs
are illustrative placeholders that I selected myself; they are not measurements, not elicited from
practitioners, and not validated. No interview and no independent expert review has yet been
completed. The study asserts nothing about whether any system is compliant with FDA requirements or
ready to manufacture drugs, and it makes no claim about the safety or quality of any product or
process. What the model can do is show which constraints determine the outcome, and that structural
result is stable across every revision of the model, including the revisions that overturned its
other conclusions.

Two of the three comments below respond to questions the Agency expressly raised. I offer them as
structural observations about where the proposed pathway is likely to bind, not as evidence about
any firm or product.

\section{Please define ``unified pharmaceutical quality system,'' because disposition authority is
the load-bearing term}

The preamble notes that ``unified pharmaceutical quality system'' and ``equivalent'' are not
currently defined in FDA regulations, and seeks comment on whether they should be defined in the
final rule and how. They should be, and the specific ambiguity worth resolving is which quality-unit
functions may be exercised centrally across distributed manufacturing units and which must remain
resident at each unit.

This is not a drafting preference. In my model, every multi-site architecture rests on an assumption
that a single quality unit can perform batch disposition for several registered locations, with only
some functions resident locally. I recorded that assumption as unsourced, because my only basis for
it is my own reading of 21 CFR 211.22(a), which assigns the quality control unit responsibility for
approving or rejecting each batch without addressing the distributed case. If that assumption is
wrong, an entire family of distributed designs collapses, not because of cost but because each unit
then carries a full resident quality function and the configuration stops being distributed in any
economically meaningful sense.

A registration pathway that permits a distributed manufacturing establishment to register as one
establishment, while leaving unresolved whether its quality system may dispose of batches as one
quality unit, gives a sponsor the administrative benefit without the operational one. A definition
of ``unified pharmaceutical quality system'' that speaks directly to disposition authority, and to
the minimum functions that must remain at each unit, would let a sponsor determine before investing
whether the pathway is usable at all.

\section{The registration timelines are not the constraint that governs supply; time to first
released batch is}

The Agency specifically requests comment on the appropriateness of the 30-day and 120-day timelines
for notifying FDA of a mobile unit relocation. On the narrow question I have no operational basis to
argue for a different number, and I do not propose one.

I would offer a different observation about those timelines. In my model the interval that decides
whether added capacity reduces a shortage is not the registration notification interval. It is the
elapsed time from the decision to add capacity to the first batch that can actually be released and
shipped, which includes qualification, process validation, and the release hold itself. Under my
illustrative inputs, sterility testing sits on the critical path of that release hold and dominates
it: chemical and analytical testing runs concurrently and its removal would shorten release by a
small fraction of the total, while the incubation period does not compress. The consequence is that
a distributed unit contributes nothing during the period before it exists as a releasing site, and
in the model that dead interval, rather than any steady-state efficiency, is what most often
determines whether the distributed configuration meets a service requirement.

The implication for this rulemaking is modest but I think real. A registration pathway that permits
a unit to be added by expedited update removes an administrative interval that is short relative to
the qualification and release interval that follows it. That is a genuine benefit and worth having.
It is unlikely, on its own, to change how quickly distributed capacity can respond to a shortage,
and I would suggest the final rule's benefit discussion distinguish the registration interval from
the time to first released batch rather than treating flexibility in the former as a proxy for
responsiveness.

\section{The preliminary economic analysis may overstate uptake if it does not separate the
capacity-short case}

My model's clearest finding, and the one that survived every correction I made, is that a network
whose saleable capacity is below demand cannot be stocked or scheduled out of the deficit. In the
runs I have done, the configurations that meet a service requirement where the incumbent plant is
capacity-short are not distributed plants that a sponsor builds. They are arrangements that buy
campaigns on lines that are already registered and already inspected, paired with positioned
finished-goods inventory. Distributed capacity performed best in my model on the product that was
\emph{not} capacity-short, which is the case where a shortage is least likely to arise.

I state this carefully, because the inputs are illustrative and the result is model behavior rather
than a fact about the industry. But if the pattern holds under real inputs, it suggests the sponsors
most likely to use this registration pathway are those already operating registered lines and
seeking flexibility across them, rather than new entrants building regional capacity to address
shortages. That distinction matters for the preliminary economic analysis, since the two populations
differ in number, in size, and in how much the administrative saving is worth to them. I would
encourage the Agency to consider whether its uptake estimate rests on the second population when the
first is the more likely user.

\section{What I am not asking for}

I am not asking FDA to endorse, review, or respond to my model, and nothing in this comment should
be read as a claim that the Agency has done so. I am not seeking any approval, designation, meeting,
or determination of any kind, and I am not requesting a response. Nobody asked me to submit this and
nobody is paying me to. My interest in the subject is disclosed in section 1 and I have not tried to
write around it: the analysis below is unfavorable to the configuration I originally set out to
build, which is the main reason I think it is worth the Agency's time.

I am a student who spent a year modeling the configuration this rule addresses, found that my own
starting hypothesis did not survive, and thought the three constraints above might be useful to the
docket. The full study, including its assumptions, its defect register, and the conclusions it has
retracted, is available on request and will be posted publicly.

\vspace{1.2ex}
\hrule
\vspace{0.8ex}
{\footnotesize Every quantitative statement above describes the behavior of a computational model
under illustrative inputs that have not been validated, elicited, or peer reviewed. None is offered
as a measurement, a regulatory position, or a statement about any firm, product, or facility.}

\end{document}
